\section{Rectas con punto flotante}
\whenuse{Proyecci\'on, distancia e intersecci\'on de rectas cuando el resultado no
necesita ser entero y se controla el error con epsilon.}
\complexity{cada operaci\'on de rectas $O(1)$ tiempo y memoria.}

\begin{lstlisting}[style=cpp]
const long double GEOMETRY_EPS = 1e-12L;

struct DPoint {
    long double x, y;
    DPoint operator+(DPoint other) const { return {x + other.x, y + other.y}; }
    DPoint operator-(DPoint other) const { return {x - other.x, y - other.y}; }
    DPoint operator*(long double scale) const { return {x * scale, y * scale}; }
    DPoint operator/(long double scale) const { return {x / scale, y / scale}; }
};

long double dot(DPoint a, DPoint b) { return a.x * b.x + a.y * b.y; }
long double cross(DPoint a, DPoint b) { return a.x * b.y - a.y * b.x; }
long double norm(DPoint a) { return sqrtl(dot(a, a)); }

optional<DPoint> line_intersection(DPoint a, DPoint b,
                                   DPoint c, DPoint d) {
    DPoint first = b - a, second = d - c;
    long double denominator = cross(first, second);
    if (fabsl(denominator) <= GEOMETRY_EPS) return nullopt;
    long double parameter = cross(c - a, second) / denominator;
    return a + first * parameter;
}

DPoint projection_on_line(DPoint point, DPoint a, DPoint b) {
    DPoint direction = b - a;
    assert(norm(direction) > GEOMETRY_EPS);
    return a + direction * (dot(point - a, direction) / dot(direction, direction));
}

long double distance_to_line(DPoint point, DPoint a, DPoint b) {
    DPoint direction = b - a;
    assert(norm(direction) > GEOMETRY_EPS);
    return fabsl(cross(direction, point - a)) / norm(direction);
}
\end{lstlisting}

\begin{lstlisting}[style=cpp]
DPoint a{0,0},b{4,0},c{2,-2},d{2,2},p{3,3};
auto intersection=line_intersection(a,b,c,d);
cout<<(double)intersection->x<<' '<<(double)intersection->y<<'\n'; // 2 0
DPoint projected=projection_on_line(p,a,b);
cout<<(double)projected.x<<' '<<(double)projected.y<<'\n'; // 3 0
cout<<(double)distance_to_line(p,a,b)<<'\n'; // 3
// nullopt en line_intersection significa rectas paralelas o coincidentes.
\end{lstlisting}

\section{C\'irculos}
\whenuse{Intersecciones recta--c\'irculo o c\'irculo--c\'irculo, tangencias y
distancias asociadas a discos.}
\complexity{cada operaci\'on de c\'irculos $O(1)$ tiempo; memoria $O(1)$ m\'as la salida.}

\begin{lstlisting}[style=cpp]
struct Circle { DPoint center; long double radius; };

// nullopt: circulos coincidentes (infinitas intersecciones).
optional<vector<DPoint>> circle_intersections(Circle first, Circle second) {
    DPoint difference = second.center - first.center;
    long double distance = norm(difference);
    if (distance <= GEOMETRY_EPS &&
        fabsl(first.radius - second.radius) <= GEOMETRY_EPS) return nullopt;
    if (distance > first.radius + second.radius + GEOMETRY_EPS ||
        distance < fabsl(first.radius - second.radius) - GEOMETRY_EPS)
        return vector<DPoint>{};

    long double along = (first.radius * first.radius
                        - second.radius * second.radius
                        + distance * distance) / (2 * distance);
    long double height_squared = first.radius * first.radius - along * along;
    if (height_squared < 0 && height_squared > -GEOMETRY_EPS) height_squared = 0;
    if (height_squared < 0) return vector<DPoint>{};

    DPoint unit = difference / distance;
    DPoint base = first.center + unit * along;
    long double height = sqrtl(height_squared);
    DPoint perpendicular{-unit.y, unit.x};
    if (height <= GEOMETRY_EPS) return vector<DPoint>{base};
    return vector<DPoint>{base + perpendicular * height,
                          base - perpendicular * height};
}
\end{lstlisting}

\begin{lstlisting}[style=cpp]
Circle first{{0,0},2},second{{2,0},2};
auto intersections=circle_intersections(first,second);
if(!intersections) cout<<"Coincidentes\n"; // Infinitos puntos comunes.
else cout<<intersections->size()<<'\n';     // 2: (1,+sqrt(3)), (1,-sqrt(3))
// size 0: separados/contenidos; size 1: tangentes; size 2: se cortan.
\end{lstlisting}

\section{Punto en pol\'igono}
\whenuse{Decidir si un punto est\'a dentro, fuera o sobre el borde de un pol\'igono
simple, incluso no convexo.}
\complexity{$O(n)$ tiempo y $O(1)$ memoria adicional.}

Devuelve 0 afuera, 1 adentro y 2 sobre el borde. Funciona para pol\'igonos
simples en cualquier orientaci\'on.

\begin{lstlisting}[style=cpp]
bool on_segment(DPoint a, DPoint b, DPoint point) {
    if (fabsl(cross(b - a, point - a)) > GEOMETRY_EPS) return false;
    return dot(point - a, point - b) <= GEOMETRY_EPS;
}

int point_in_polygon(const vector<DPoint> &polygon, DPoint point) {
    bool inside = false;
    int n = polygon.size();
    for (int i = 0; i < n; ++i) {
        DPoint a = polygon[i], b = polygon[(i + 1) % n];
        if (on_segment(a, b, point)) return 2;
        bool crosses = (a.y > point.y) != (b.y > point.y);
        if (crosses) {
            long double intersection_x = a.x
                + (b.x - a.x) * (point.y - a.y) / (b.y - a.y);
            if (point.x < intersection_x) inside = !inside;
        }
    }
    return inside ? 1 : 0;
}
\end{lstlisting}

\begin{lstlisting}[style=cpp]
vector<DPoint> polygon={{0,0},{4,0},{4,4},{0,4}};
cout<<point_in_polygon(polygon,{2,2})<<'\n'; // 1: adentro
cout<<point_in_polygon(polygon,{0,2})<<'\n'; // 2: sobre el borde
cout<<point_in_polygon(polygon,{5,2})<<'\n'; // 0: afuera
\end{lstlisting}

\section{Rotating calipers: di\'ametro convexo}
\whenuse{Obtener el par de puntos m\'as alejados de un pol\'igono convexo; aplicar
primero convex hull si la entrada es un conjunto desordenado.}
\complexity{$O(n)$ tiempo y $O(1)$ memoria adicional.}

El pol\'igono debe ser convexo, sin puntos repetidos y con sus v\'ertices en
orden sobre el borde, horario o antihorario. Se detecta el sentido una sola
vez y se elige la comparaci\'on correspondiente, sin copiar ni invertir el
pol\'igono. Eliminar v\'ertices colineales intermedios con \texttt{convex\_hull}
antes de usar calipers. Si todos son colineales, se toman los extremos.

\begin{lstlisting}[style=cpp]
template<bool Clockwise>
i128 convex_diameter_oriented(const vector<Point> &polygon) {
    int n = polygon.size();
    i128 answer = 0;
    int opposite = 1;
    for (int i = 0; i < n; ++i) {
        int next = (i + 1) % n;
        while (true) {
            int candidate = (opposite + 1) % n;
            i128 current_area = cross(polygon[next] - polygon[i],
                                      polygon[opposite] - polygon[i]);
            i128 next_area = cross(polygon[next] - polygon[i],
                                   polygon[candidate] - polygon[i]);
            if constexpr (Clockwise) {
                if (next_area >= current_area) break;
            } else {
                if (next_area <= current_area) break;
            }
            opposite = candidate;
        }
        answer = max(answer, squared_distance(polygon[i], polygon[opposite]));
        answer = max(answer, squared_distance(polygon[next], polygon[opposite]));
    }
    return answer;
}

i128 convex_diameter_squared(const vector<Point> &polygon) {
    int n = polygon.size();
    if (n <= 1) return 0;
    if (n == 2) return squared_distance(polygon[0], polygon[1]);
    for (int i = 0; i < n; ++i) {
        i128 turn = cross(polygon[i], polygon[(i + 1) % n],
                          polygon[(i + 2) % n]);
        if (turn > 0) return convex_diameter_oriented<false>(polygon);
        if (turn < 0) return convex_diameter_oriented<true>(polygon);
    }
    // Todos colineales: los extremos lexicograficos son los mas lejanos.
    auto [first, last] = minmax_element(polygon.begin(), polygon.end());
    return squared_distance(*first, *last);
}
\end{lstlisting}

\begin{lstlisting}[style=cpp]
vector<Point> points={{0,0},{2,0},{2,2},{0,2},{1,1}};
vector<Point> hull=convex_hull(points);
i128 diameter_squared=convex_diameter_squared(hull);
cout<<(ll)diameter_squared<<'\n'; // 8: distancia entre esquinas opuestas
reverse(hull.begin(),hull.end()); // El mismo poligono en sentido horario.
cout<<(ll)convex_diameter_squared(hull)<<'\n'; // Tambien 8
// La distancia real es sqrt(8); se conserva al cuadrado para evitar decimales.
\end{lstlisting}
